\PassOptionsToPackage{unicode}{hyperref}
\PassOptionsToPackage{hyphens}{url}
\documentclass[11pt,a4paper]{article}

% ── Codificação e língua ──────────────────────────────────────────────────────
\usepackage[T1]{fontenc}
\usepackage[utf8]{inputenc}
\usepackage[brazil]{babel}

% ── Tipografia ────────────────────────────────────────────────────────────────
\usepackage{lmodern}
\usepackage{microtype}

% ── Cores ─────────────────────────────────────────────────────────────────────
\usepackage{xcolor}
\definecolor{javaOrange}{HTML}{E76F00}
\definecolor{javaDeep}{HTML}{BF5700}
\definecolor{javaBlue}{HTML}{1565C0}
\definecolor{javaTeal}{HTML}{00695C}
\definecolor{javaAmber}{HTML}{B45309}
\definecolor{javaInk}{HTML}{2B2140}
\definecolor{javaCodeBg}{HTML}{FFF8F0}
\definecolor{javaRule}{HTML}{FFD7B5}
\definecolor{javaShade}{HTML}{FFF3E8}

% ── Caixas e código ───────────────────────────────────────────────────────────
\usepackage{tcolorbox}
\tcbuselibrary{skins,breakable,listings}
\usepackage{listings}
\usepackage{fvextra}

\lstdefinestyle{javastyle}{
  language=Java,
  basicstyle=\ttfamily\small,
  keywordstyle=\color{javaBlue}\bfseries,
  stringstyle=\color{javaDeep},
  commentstyle=\color{javaTeal}\itshape,
  numberstyle=\tiny\color{gray},
  breaklines=true,
  breakatwhitespace=false,
  showstringspaces=false,
  tabsize=4,
}

\newtcblisting{javacode}{
  listing only,
  listing style=javastyle,
  breakable,
  enhanced,
  colback=javaCodeBg,
  colframe=javaDeep,
  leftrule=4pt,
  rightrule=0.4pt,
  toprule=0.4pt,
  bottomrule=0.4pt,
  arc=2pt,
  fontupper=\small,
  top=4pt, bottom=4pt, left=6pt, right=6pt,
}

\newtcbox{\inlinecode}{
  on line,
  colback=javaCodeBg,
  colframe=javaRule,
  boxrule=0.5pt,
  arc=2pt,
  fontupper=\ttfamily\small\color{javaDeep},
  boxsep=1pt, left=2pt, right=2pt, top=1pt, bottom=1pt,
}

% ── Layout ────────────────────────────────────────────────────────────────────
\usepackage[margin=2.4cm, top=2.8cm, bottom=2.8cm]{geometry}

% ── Cabeçalho / rodapé ────────────────────────────────────────────────────────
\usepackage{fancyhdr}
\pagestyle{fancy}
\fancyhf{}
\renewcommand{\headrulewidth}{0.4pt}
\renewcommand{\headrule}{\hbox to\headwidth{\color{javaRule}\leaders\hrule height \headrulewidth\hfill}}
\fancyhead[L]{\small\sffamily\color{javaDeep} Java Programming MOOC}
\fancyhead[R]{\small\sffamily\color{javaDeep} Resoluções · Lição 4}
\fancyfoot[C]{\small\sffamily\color{gray}\thepage}

% ── Seções ────────────────────────────────────────────────────────────────────
\usepackage{titlesec}
\titleformat{\section}
  {\sffamily\Large\bfseries\color{javaOrange}}
  {\thesection}{0.7em}{}
  [\vspace{2pt}{\color{javaRule}\titlerule[1.2pt]}]
\titleformat{\subsection}
  {\sffamily\large\bfseries\color{javaDeep}}
  {\thesubsection}{0.6em}{}
\titleformat{\subsubsection}
  {\sffamily\normalsize\bfseries\color{javaTeal}}
  {\thesubsubsection}{0.5em}{}

% ── Listas ────────────────────────────────────────────────────────────────────
\usepackage{enumitem}
\setlist[itemize]{itemsep=2pt, topsep=4pt, parsep=0pt}
\setlist[enumerate]{itemsep=2pt, topsep=4pt, parsep=0pt}

% ── Tabelas ───────────────────────────────────────────────────────────────────
\usepackage{booktabs}
\usepackage{array}
\usepackage{colortbl}
\usepackage{longtable}

% ── Hiperlinks ────────────────────────────────────────────────────────────────
\usepackage{hyperref}
\hypersetup{
  colorlinks=true,
  linkcolor=javaBlue,
  urlcolor=javaBlue,
  citecolor=javaBlue,
  pdftitle={Programação em Java --- Resoluções --- Lição 4},
  pdfauthor={Java Programming MOOC · Universidade de Helsinque},
  pdflang={pt-BR},
}

% ── Caixa de destaque (dica/atenção) ─────────────────────────────────────────
\newtcolorbox{destaque}[1][Observação]{
  enhanced, breakable,
  colback=javaShade,
  colframe=javaOrange,
  leftrule=4pt, rightrule=0.4pt, toprule=0.4pt, bottomrule=0.4pt,
  arc=2pt,
  fonttitle=\sffamily\bfseries\color{javaOrange},
  title=#1,
  top=4pt, bottom=4pt,
}

% ── Título personalizado ──────────────────────────────────────────────────────
\makeatletter
\newcommand{\subtitle}[1]{\gdef\@subtitle{#1}}
\renewcommand{\maketitle}{%
  \begin{center}
    {\color{javaRule}\rule{\linewidth}{1pt}}\\[6pt]
    {\Huge\sffamily\bfseries\color{javaOrange} \@title}\\[6pt]
    \ifx\@subtitle\undefined\else
      {\Large\sffamily\color{javaDeep} \@subtitle}\\[4pt]
    \fi
    {\small\sffamily\color{gray} \@author}\\[2pt]
    {\small\sffamily\color{gray} \@date}\\[6pt]
    {\color{javaRule}\rule{\linewidth}{1pt}}
  \end{center}
  \vspace{1em}
}
\makeatother

% ─────────────────────────────────────────────────────────────────────────────
\title{Programação em Java — Resoluções}
\subtitle{Lição 4}
\author{Java Programming MOOC · Universidade de Helsinque — tradução para o português}
\date{2026}
% ─────────────────────────────────────────────────────────────────────────────

\begin{document}
\maketitle
\tableofcontents
\vspace{1.5em}

% ─────────────────────────────────────────────────────────────────────────────
\section{Exercício 1 --- Retângulo}

\begin{javacode}
public class Retangulo {
    private double largura;
    private double altura;

    public Retangulo(double largura, double altura) {
        this.largura = largura;
        this.altura = altura;
    }

    public double getArea() {
        return largura * altura;
    }

    public double getPerimetro() {
        return 2 * (largura + altura);
    }

    public boolean ehQuadrado() {
        return largura == altura;
    }

    @Override
    public String toString() {
        return "Retangulo{largura=" + largura + ", altura=" + altura +
               ", area=" + getArea() + ", perimetro=" + getPerimetro() + "}";
    }

    public static void main(String[] args) {
        Retangulo r1 = new Retangulo(4, 4);
        Retangulo r2 = new Retangulo(3, 5);

        System.out.println(r1);
        System.out.println("É quadrado? " + r1.ehQuadrado());

        System.out.println(r2);
        System.out.println("É quadrado? " + r2.ehQuadrado());
    }
}
\end{javacode}

% ─────────────────────────────────────────────────────────────────────────────
\section{Exercício 2 --- Aluno e notas}

\begin{javacode}
import java.util.ArrayList;

public class Aluno {
    private String nome;
    private ArrayList<Integer> notas;

    public Aluno(String nome) {
        this.nome = nome;
        this.notas = new ArrayList<>();
    }

    public void adicionarNota(int nota) {
        if (nota < 0 || nota > 10) {
            throw new IllegalArgumentException("A nota deve estar entre 0 e 10.");
        }
        notas.add(nota);
    }

    public double calcularMedia() {
        if (notas.isEmpty()) {
            return 0;
        }
        int soma = 0;
        for (int nota : notas) {
            soma += nota;
        }
        return (double) soma / notas.size();
    }

    public boolean aprovado() {
        return calcularMedia() >= 6;
    }

    @Override
    public String toString() {
        String situacao = aprovado() ? "Aprovado" : "Reprovado";
        return nome + " - média: " + calcularMedia() + " - " + situacao;
    }

    public static void main(String[] args) {
        Aluno aluno = new Aluno("Pedro");
        aluno.adicionarNota(8);
        aluno.adicionarNota(7);
        aluno.adicionarNota(5);

        System.out.println(aluno);
    }
}
\end{javacode}

% ─────────────────────────────────────────────────────────────────────────────
\section{Exercício 3 --- Estacionamento}

\begin{javacode}
public class Veiculo {
    private String placa;
    private String marca;

    public Veiculo(String placa, String marca) {
        this.placa = placa;
        this.marca = marca;
    }

    public String getPlaca() {
        return placa;
    }

    public String getMarca() {
        return marca;
    }

    @Override
    public String toString() {
        return marca + " - " + placa;
    }
}
\end{javacode}

\begin{javacode}
import java.util.ArrayList;

public class Estacionamento {
    private int capacidade;
    private ArrayList<Veiculo> veiculos;

    public Estacionamento(int capacidade) {
        this.capacidade = capacidade;
        this.veiculos = new ArrayList<>();
    }

    public boolean entrar(Veiculo v) {
        if (veiculos.size() >= capacidade) {
            System.out.println("Estacionamento lotado. Não é possível entrar.");
            return false;
        }
        veiculos.add(v);
        return true;
    }

    public boolean sair(String placa) {
        for (Veiculo v : veiculos) {
            if (v.getPlaca().equals(placa)) {
                veiculos.remove(v);
                return true;
            }
        }
        return false;
    }

    public void listarVeiculos() {
        System.out.println("Veículos no estacionamento (" + veiculos.size() + "/" + capacidade + "):");
        for (Veiculo v : veiculos) {
            System.out.println("- " + v);
        }
    }

    public static void main(String[] args) {
        Estacionamento estacionamento = new Estacionamento(2);

        estacionamento.entrar(new Veiculo("ABC1234", "Fiat"));
        estacionamento.entrar(new Veiculo("XYZ9876", "Toyota"));
        estacionamento.entrar(new Veiculo("QWE4567", "Ford")); // lotado

        estacionamento.listarVeiculos();

        estacionamento.sair("ABC1234");
        estacionamento.listarVeiculos();
    }
}
\end{javacode}

% ─────────────────────────────────────────────────────────────────────────────
\section{Exercício 4 --- Leitura de arquivo}

Primeiro, criamos o arquivo \inlinecode{produtos.txt} com linhas no formato \inlinecode{nome,preco,estoque}:

\begin{verbatim}
Caneta,2.50,100
Caderno,15.90,40
Mochila,89.90,12
Estojo,22.30,25
\end{verbatim}

\begin{javacode}
public class Produto {
    private String nome;
    private double preco;
    private int estoque;

    public Produto(String nome, double preco, int estoque) {
        this.nome = nome;
        this.preco = preco;
        this.estoque = estoque;
    }

    public double getValorTotal() {
        return preco * estoque;
    }

    @Override
    public String toString() {
        return nome + " - R$" + preco + " - estoque: " + estoque;
    }
}
\end{javacode}

\begin{javacode}
import java.io.BufferedReader;
import java.io.FileReader;
import java.io.IOException;
import java.util.ArrayList;

public class LeitorDeProdutos {
    public static void main(String[] args) {
        ArrayList<Produto> produtos = new ArrayList<>();

        try (BufferedReader leitor = new BufferedReader(new FileReader("produtos.txt"))) {
            String linha;
            while ((linha = leitor.readLine()) != null) {
                String[] partes = linha.split(",");

                String nome = partes[0];
                double preco = Double.valueOf(partes[1]);
                int estoque = Integer.valueOf(partes[2]);

                produtos.add(new Produto(nome, preco, estoque));
            }
        } catch (IOException e) {
            System.out.println("Erro ao ler o arquivo: " + e.getMessage());
            return;
        }

        double valorTotal = 0;

        System.out.println("Produtos cadastrados:");
        for (Produto produto : produtos) {
            System.out.println(produto);
            valorTotal += produto.getValorTotal();
        }

        System.out.println("\nValor total em estoque: R$" + valorTotal);
    }
}
\end{javacode}

\begin{destaque}
O uso do \inlinecode{try-with-resources} garante que o arquivo seja fechado automaticamente mesmo se ocorrer um erro durante a leitura.
\end{destaque}

\end{document}
